\subsection{The actual inverse-parameter kernel estimate}
\label{actual-off-diagonal}

Keep the original place, tower, forms, finite grids and positive
normalizer $n$ of the preceding Fourier theorem.  Write
$N_{j,w}(a,b)$ and $E_{j,w}(a)$ for the same normalized
kernel and its diagonal.  The estimate below concerns every
original admissible word, rather than only words assigned to
successive interpolation blocks.

Put $N_*=20(r+1)$, $I=\beta-\alpha\ge6N_*^3+6N_*$ and
$\eta=28N_*/I$.  Define
\[
 A_1=\max_{y\in\mathcal Y}
       \sqrt{\frac{n(y)n(y+1)}{q_0(y)}},\qquad
 \Lambda_{j,s}^{\rm arg}=e^{10/49}A_1\eta^{2j-s}
       \quad(1\le s\le j\le r).
\]
For $r\ge1$ let $\Lambda^{\rm arg}$ be the maximum of
these finitely many positive numbers.  For $r=0$ the
statement below has no instances, so any positive value
may be used.  A common scale for this theorem and the
preceding Fourier theorem is the maximum of their two
scales, enlarged to one if the consumer requires it.

\begin{theorem}\label{off-diagonal-actual}
Let $F$ be an actual smooth compact coefficient with positive
Gram cutoff.  For $1\le s\le j\le r$, let $h$ be any
relevant head of size $j-s$, and let $t,t'\ne0$.
For every original admissible word $w$ on $\mathcal Y$,
every $K\ge0$ with $|w|+K\le N_*$, there are two original
extensions $w_1,w_2$, each obtained by prepending at most
$K$ letters and having length at most $N_*$, such that
\begin{align}\label{off-diagonal-bound}
 &\left|N_{j,w}\bigl(h\mathbin{\|}(s,t),
                           h\mathbin{\|}(s,t')\bigr)\right|\notag\\
 &\quad\le
 m\bigl(\Lambda^{\rm arg},K,|t^{-1}-{t'}^{-1}|\bigr)
 \frac{E_{j,w_1}(h\mathbin{\|}(s,t))+
       E_{j,w_2}(h\mathbin{\|}(s,t'))}{2}.
\end{align}
Here $m(\Lambda,K,0)=1$ and
$m(\Lambda,K,u)=\min(1,(\Lambda/u)^K)$ for $u>0$.
This supplies the original
\texttt{KernelArgumentHypothesisScalar} on these grids.
The proof also supplies the original \texttt{LiftDegAll}
and \texttt{DegLift}, with their full compact-coefficient,
head, word and interpolation quantifiers.
\end{theorem}

\begin{proof}
We first prove the missing identity for the prescribed
operators, before applying interpolation.
Use the paired basis $e_0,\ldots,e_{r-1},w_0,\ldots,w_k,
f_{r-1},\ldots,f_0$ and the actual complexified right
derivatives $E_{ab}$.  They satisfy
\[
 [E_{ab},E_{cd}]=\delta_{bc}E_{ad}-\delta_{da}E_{cb}.
\]
For ordered injective lists $a,b$ of length $m$, let
$Q_{a,b}(u)$ be the ordered column minor whose factors
are $T_{ab}(u-\ell)$, with
$T_{ab}(v)=v\delta_{ab}+E_{ab}$ and $0\le\ell<m$.
The rightmost factor acts first.  Empty minors are one.
The last-column and last-row formulas are
\begin{align}
 Q_{a,b}(u)
 &=\sum_{\ell=0}^{m-1}(-1)^{m-1+\ell}
 Q_{a\setminus a_\ell,\,b\setminus b_{m-1}}(u)
                  T_{a_\ell b_{m-1}}(u-m+1),
                                      \label{off-diagonal-col}\\
 Q_{a,b}(u)
 &=\sum_{\ell=0}^{m-1}(-1)^{m-1+\ell}
 Q_{a\setminus a_{m-1},\,b\setminus b_\ell}(u-1)
                  T_{a_{m-1}b_\ell}(u).
                                      \label{off-diagonal-row}
\end{align}
The first follows by grouping permutations according to
their last entry.  For the second, the Lie relation gives
the two-factor identity
\[
 T_{ab}(u)T_{cd}(u-1)-T_{cb}(u)T_{ad}(u-1)
 =T_{ab}(u-1)T_{cd}(u)-T_{ad}(u-1)T_{cb}(u).
\]
Apply this identity successively to adjacent positions in
the alternating sum defining the minor.  The terms with
the two exchanged row indices have opposite signs;
the displayed identity replaces their column order by
row order and moves the spectral shift with the position.
Induction on the number of positions gives the ordered
row formula, whose last-row expansion is
\eqref{off-diagonal-row}.  This argument retains the
$u-1$ in that formula; treating the generators as
commuting would give the wrong dual peeling identity
\cite{off-diagonal-source}.

Let $\mathcal R H$ be the original right average
\[
 (\mathcal R H)(g)=\int_{K'\times N'}
 \overline{\psi_{z',-2i}(n)}\,H(gnk)
                         \pi'(k)^{-1}\,dk\,dn.
\]
Both measures are the original tower measures.  Left
invariance of $dn$, without a change of normalization,
gives $(\mathcal R H)(gn)=\psi_{z',-2i}(n)(\mathcal R H)(g)$.
Differentiating the original character and complexifying
its two real tangent directions gives
\begin{equation}\label{off-diagonal-character}
 E_{ab}\mathcal R H=\chi_{ab}\mathcal R H,
 \qquad
 \chi_{e_a,e_b}=\mathbf1_{b=a+1},\qquad
 \chi_{f_a,f_b}=\mathbf1_{a=b+1}.
\end{equation}
These equations are used only in the unipotent root
directions.  In particular, the second formula uses the
displayed reversed ambient order on the $f_a$'s.
The middle-block character values retain the original
$z'$; they are not replaced, and are not needed for
the two chain peelings below.

For a fixed $g$, the functional
$\mathcal L(X)=(X\mathcal R H)(g)$ consequently satisfies
$\mathcal L(XE_{ab})=\chi_{ab}\mathcal L(X)$ for every
operator word $X$, whenever $E_{ab}$ is one of these
root directions.  Equivalently
it kills the left ideal generated by $E_{ab}-\chi_{ab}$.
The generators occur on the \emph{right} in this rule.

Put $i=j-s$.  Apply \eqref{off-diagonal-col} to the minor
with rows $e_0,\ldots,e_{j-2}$ and columns
$(e_0,\ldots,e_{j-1})\setminus e_i$.
In its last column, every factor is off diagonal.
Equation \eqref{off-diagonal-character} kills all terms
except the row $e_{j-2}$, whose character value and
cofactor sign are both one.  Repeat until reaching
the principal $i\times i$ minor.  Thus, modulo that
left ideal,
\[
 Q_{(e_0,\ldots,e_{j-2}),\,
           (e_0,\ldots,e_{j-1})\setminus e_i}(u)
 \equiv Q_{(e_0,\ldots,e_{i-1}),\,
                        (e_0,\ldots,e_{i-1})}(u).
\]
For the dual minor use \eqref{off-diagonal-row}, starting
with rows $(f_0,\ldots,f_{j-1})\setminus f_i$,
columns $f_0,\ldots,f_{j-2}$, and parameter $v+j-2$.
Only the last column survives, again with coefficient one.
Each peeling lowers the parameter by one, so the final
principal $i\times i$ minor has parameter $v+i-1$.
If $s=1$ the respective omitted index is already last
and neither peeling is needed.  If $i=0$ the final
minor is the empty minor.  These cover both edge cases.

Now use the actual normal representative of
$a=h\mathbin{\|}(s,t)$.  Its first $j$ paired rows form
$\operatorname{diag}(H_h,tJ_s)$, where $J_s$ is the
reversing permutation matrix.  Write
$\epsilon_s=\det J_s=(-1)^{s(s-1)/2}$.
The original principal minors are
\[
 D_i=\det H_h,\qquad D_j=\det H_h\,t^s\epsilon_s.
\]
In the contracted minor of size $j-1$, the only nonzero
column set omits column $i$.  Its coefficient is
$\det H_h\,t^{s-1}\epsilon_{s-1}$.
The dual contraction has its complex conjugate and the
reversed $f$ order.  The preceding peelings, followed by
division by $D_j$ or $\overline{D_j}$, therefore give
\begin{equation}\label{off-diagonal-letter-reduction}
 \begin{split}
 L^+_{j,1}(y)\mathcal R H
    &=(-1)^{s-1}t^{-1}L^+_{i,0}(y)\mathcal R H,\\
 L^-_{j,1}(y)\mathcal R H
    &=(-1)^{s-1}\bar t^{-1}L^-_{i,0}(y)\mathcal R H.
 \end{split}
\end{equation}
These equalities hold at that representative and at
every translate by the original left flag stabilizer.
Indeed, its action fixes the paired rows and their
minors, and the contracted differential operators have
the already proved left covariance.  The spectral
parameters here are the actual $y+(2r+k)/2$ and
$y-(2r+k)/2$; the dual size-dependent shift in the
minor was explicitly retained in its peeling.

We justify applying this calculation to the actual
slices.  Every old-word component of $F$ has compact
support and is smooth on $D_j\ne0$.  Localizing in
$D_j$ around the fixed nonzero value $D_j(a)$ gives
a globally smooth compact function without changing
any value or derivative on the integration locus.
The groups in the right average fix this minor.
For $g$ in a compact neighborhood, the portion of
$N'K'$ meeting the support of $H(gnk)$ is compact:
$nk$ lies in the compact set $g^{-1}\operatorname{supp}H$,
$k$ lies in a compact group, and $N'$ is closed.
This proves smoothness of the average and permits
differentiation under its integral.  The exact right
covariance of the contracted minors cancels the
conjugation of the right derivatives.  Hence
\[
       \mathcal R(L^\pm_{j,e}(y)H)
                    =L^\pm_{j,e}(y)\mathcal R H
                    \quad(e=0,1).
\]
There is no averaging constant in this equality.

The proved head quotient factorization expresses the
Jacquet transform of a slice as the original left
quotient integral of $\mathcal R H$ at precisely the
translates used in \eqref{off-diagonal-letter-reduction}.
All integrals are absolutely convergent by the actual
compact-family properness theorem; its measure
factorization is the original normalized tower's
factorization.  Since $L^\pm_{i,0}(y)$ is a polynomial
of degree $i$ with leading coefficient the identity,
\eqref{off-diagonal-letter-reduction} proves:
after this Jacquet transform, the new bit-one slice
has degree at most $j-s$, with leading coefficient
$(-1)^{s-1}t^{-1}$ times the old slice for plus and
$(-1)^{s-1}\bar t^{-1}$ times the old slice for minus.
Any fixed real shift of $y$ leaves these leading
coefficients unchanged.  Its coefficients may be
integrated separately, since finite Vandermonde
inversion expresses each as a finite linear
combination of the absolutely integrable evaluations.
This is the full original \texttt{LiftDegAll}.

The degree decrease has been proved \emph{after the
Jacquet transform}.  It need not hold for the slice
function itself.  To pass to the original residue
form, take the polynomial coefficients supplied by
the compact slice theorem.  Uniqueness of polynomial
coefficients in their Jacquet transforms allows us
to set coefficients above degree $j-s$ to zero and
replace the top one by the stated multiple of the
old slice.  Each difference has zero Jacquet
transform.  The proved exact Jacquet Gram identity,
together with the intertwining and preservation of
the old residue operators, makes each such difference
invisible in the original form and in every mixed
contraction of that form.  This uses the actual
compact coefficient class on both sides; it asserts
no positivity on a larger radial space.

Let $a=h\mathbin{\|}(s,t)$ and
$a'=h'\mathbin{\|}(s,t')$, allowing different preceding
heads of the same size for this algebraic step.
The fresh mixed contraction is
$\beta(x_1,y_0)-\beta(x_0,y_1)$, with the plus
insertion at $y$ and the minus insertion at $y+1$.
The bit-zero degrees are $j$, with the old vectors
as leading coefficients.  The Jacquet-reduced
bit-one degrees just proved are $j-s$.
Sesquilinearity, linear in the first argument,
therefore gives degree at most $d=2j-s$ and leading
coefficient
\begin{equation}\label{off-diagonal-leading}
        (-1)^{s-1}(t^{-1}-{t'}^{-1})
                       \beta(x_{w,a},x_{w,a'}).
\end{equation}
In particular, conjugating the negative leading
coefficient removes the bar on $t'$.
Iterating this identity tail first through any
finite list of new parameters and extracting each
leading coefficient on any injective real grid
gives the exact source \texttt{DegLift} identity.
No cut or separation is needed for that identity
beyond the actual operators' definitions.  The old
residue matrices enter only as fixed coefficients;
the same calculation also applies where the source
assigns a zero value to a scalar inverse at zero.
No cancellation of such an inverse is used here.
The nonsingular cut and separation conditions enter
the positive estimate next.

Return to $h'=h$ and put $\delta=t^{-1}-{t'}^{-1}$.
At most $|w|$ long blocks contain old underlying
nodes.  Choose $K$ distinct unoccupied long blocks;
the bound $|w|+K\le N_*$ guarantees their existence.
Since $1\le d=2j-s\le2r-1$, each chosen block
contains its already constructed degree-$d$
interpolation grid, with absolute coefficient sum
at most $\eta^d$.  Every simultaneous plus or
minus extension remains in the original cut,
and all new--old and new--new shifted distances
exceed two.  This choice works for every source
admissible word, including repeated occupancy of
a single old block.

Apply the proved actual mixed determinant inequality
to the iterated identity \eqref{off-diagonal-leading}.
For each new node retain the full correction
\[
 q_w(y)=q_0(y)
 \prod_{(+,u)\in w}\frac{y-u+1}{y-u}
 \prod_{(-,u)\in w}\frac{u-y}{u-y-1}.
\]
The original separated-grid estimate gives
$\sqrt{n(y)n(y+1)/q_w(y)}\le e^{10/49}A_1$.
Writing $a_*=n\mathrm{Pt}(a)>0$ and
$a'_* =n\mathrm{Pt}(a')>0$, the factor
$\sqrt{a_*a'_*}/S_w$ cancels exactly against
\[
 S_{w_+}=S_w\prod_\ell n(y_\ell),\qquad
 S_{w_-}=S_w\prod_\ell n(y_\ell+1)
\]
and the diagonal point normalizers.  Thus
\[
 |\delta|^K|N_{j,w}(a,a')|
 \le(\Lambda_{j,s}^{\rm arg})^K
   \sum_{\boldsymbol\nu}\theta_{\boldsymbol\nu}
       \sqrt{E_{j,w_+(\boldsymbol\nu)}(a)
             E_{j,w_-(\boldsymbol\nu)}(a')},
 \qquad
 \theta_{\boldsymbol\nu}\ge0,\quad
 \sum_{\boldsymbol\nu}\theta_{\boldsymbol\nu}=1.
\]
Positivity also gives
$|N_{j,w}(a,a')|\le\sqrt{E_{j,w}(a)E_{j,w}(a')}$.
Let $M_1,M_2$ be the separate maxima of the two
diagonals over $w$ and all the finitely many
extensions used above.  Each is attained by an
allowed $w_1,w_2$.  Both geometric means are bounded
by $\sqrt{M_1M_2}\le(M_1+M_2)/2$.
For $\delta\ne0$ take the stronger of the two
estimates; for $\delta=0$ use positivity alone.
For $K=0$ take $w_1=w_2=w$ and empty products.
This proves \eqref{off-diagonal-bound} with the
original extension quantifiers.
\end{proof}

\paragraph{Scope of the advance.}
The source writes $j=b+s-1$ with $b\ge1$;
the preceding head has size $b-1=j-s$, exactly
the theorem's range.  In particular, this consumer
does not ask for a terminal $j=r+1$ instance.
The spectral majorant from the preceding Fourier
theorem gives the explicit arithmetic bound
\[
 \Lambda_{j,s}^{\rm arg}
 \le e^{10/49}g^{-2r-2}\kappa^{2r+1}
      80t_v\sqrt{k}\,(28N_*/I)^{2j-s}
\]
under its unchanged binary-spectrum and interval
hypotheses.  All dimensions, shifts, point weights
and word corrections remain present.  Passage to
the actual decaying eigenfunctions, the block-growth
limit and the uniformly controlled top-lift supplier
are still separate analytic inputs.  This local
result alone supplies no global arithmetic Gram
estimate and completes no whole proof-map node.
